% chapters/02/aktivitas.tex
\bagian{Aktivitas Telaah Terbimbing}

Bab ini bersifat teori, sehingga aktivitasnya bukan menulis program,
melainkan menelaah permainan yang sudah ada dan menghitung besaran yang
relevan. Kerjakan ketiga telaah berikut secara berurutan; hasilnya akan
dipakai kembali pada praktikum Bab~III dan seterusnya.

\begin{aktivitasbox}
\textbf{Telaah 1: Membedah satu permainan.} Pilih satu permainan sederhana
yang Anda kenal baik, boleh permainan klasik maupun permainan bergerak.
Isilah tabel berikut berdasarkan pengamatan dan pengetahuan Anda tentang
permainan tersebut.
\end{aktivitasbox}

\begin{table}[htbp]
\centering
\caption{Lembar kerja telaah permainan}
\label{tab:lembar-telaah}
{\setlength{\tabcolsep}{4pt}\footnotesize%
\begin{tabularx}{\textwidth}{@{}L{4.5cm} Y@{}}
\toprule
\headrow \textbf{Aspek yang ditelaah} & \textbf{Jawaban Anda} \\
\midrule
Nama permainan dan alasan pemilihannya & \\[1.4em]
\addlinespace
Pemain: siapa yang mengambil keputusan, dan berapa orang? & \\[1.4em]
\addlinespace
Aturan: tindakan apa yang boleh dan tidak boleh dilakukan? & \\[2.2em]
\addlinespace
Konflik: apa yang membuat tujuan tidak dapat dicapai langsung? & \\[1.8em]
\addlinespace
Keluaran terukur: bagaimana keberhasilan dinyatakan? & \\[1.4em]
\addlinespace
Subsistem yang tampak: masukan, penyajian, audio, dan penyimpanan data &
\\[1.8em]
\addlinespace
Perkiraan laju bingkai dan alasan perkiraan tersebut & \\[1.4em]
\bottomrule
\end{tabularx}}
\end{table}

\begin{aktivitasbox}
\textbf{Telaah 2: Menghitung anggaran waktu satu bingkai.} Sebuah permainan
dijalankan pada tiga laju bingkai yang berbeda. Hitunglah lama satu bingkai
pada masing-masing laju, lalu tuliskan dua pekerjaan yang menurut Anda paling
banyak memakan waktu pada tahap simulasi dan tahap penyajian.
\end{aktivitasbox}

\begin{aktivitasbox}
\textbf{Telaah 3: Mengamati perilaku NPC.} Jalankan satu permainan yang
memiliki musuh, lalu amati perilaku musuh tersebut selama sekitar lima menit.
Catat keadaan yang Anda kenali (berdiam, berpatroli, mengejar), peristiwa yang
memicu perpindahan antar-keadaan, dan apa yang terjadi ketika Anda bersembunyi
di balik rintangan. Simpan catatan ini, karena akan Anda pakai pada Bab~VIII
dan Bab~IX.
\end{aktivitasbox}

\begin{catatabox}
Ketiga telaah tersebut tidak menuntut pemasangan perkakas apa pun. Apabila
Anda belum mengenal permainan yang akan ditelaah, pilihlah permainan sederhana
yang tersedia bebas, misalnya permainan arkade klasik yang dapat dimainkan di
peramban. Yang dinilai pada aktivitas ini adalah ketelitian pengamatan, bukan
penguasaan permainannya.
\end{catatabox}
